\section{Algorithm-Level Implementation Details}
\label{section_implementation}

    This section details the algorithm-level implementation choices that keep the online optimization lightweight: execution-time prediction, priority-value mapping, performance-metric construction, and hyperparameters.

    \subsection{Predicting Execution-Time Distributions}
    \label{section: predict_ET_exp}
    Due to the highly limited computation resources available for online optimization, we use a simple heuristic that assumes the robot's working environment will not change significantly within each short time interval.
    Let $T^W$ denote the length of such an interval where the environment vector $\textbf{E}$ changes by at most $\Theta^W$:
    \begin{equation}
        \|\textbf{E}_{t_0+T^W} - \textbf{E}_{t_0} \| \leq \Theta^W.
    \end{equation}
    Under a smoothly changing environment, $\Theta^W \approx 0$ for a sufficiently small $T^W$, so the execution-time distribution of a task $\tau_i$ over $[t_0 + T^W/2, t_0 + T^W]$ can be predicted from measurements over $[t_0, t_0 + T^W/2]$.
    In our experiments, we take $\Theta^W \approx 0$ when $T^W$ equals $\tau_i$'s release period $T_i$, yielding multiple measurements per window from which a distribution is established.
    To further reduce computation, we assume each task's execution time follows a normal distribution characterized only by its mean, variance, minimum, and maximum, reducing the overall cost to linear complexity.
    We set $T^W$ equal to the scheduler's running period for simplicity.
    These assumptions may not hold when the environment changes abruptly, leading to inaccurate predictions; however, as discussed in Section~\ref{section: relax_smooth_assumption}, such error is the necessary cost of low computation and is automatically corrected as long as the environment does not remain persistently unsmooth.


    \subsection{Mapping Priority Assignments to Priority Values}
    \label{section_pa_values}
    While the priority assignment algorithm determines the relative priority order of tasks, Linux requires specific \textit{priority values} (real-time processes range from 0 to 99, larger is higher) to implement these assignments.
    Theoretically, any mapping that adheres to the priority assignment order is valid.
    However, changing priority values involves numerous low-level OS operations, and this overhead grows for multi-threaded tasks whose threads must update simultaneously, so we minimize the number of value changes.
    We select the task with the largest number of threads as the \textit{reference task} and assign it a fixed priority value; all other tasks' values are computed relative to it, following the priority assignment order.
    This minimizes changes to the reference task's priority value; if another task's thread count surpasses the reference task's at runtime, that task is dynamically designated as the new reference.
    The strategy generalizes to multi-core systems under partitioned scheduling, where a fixed priority value can be locked for one task per core.

    \begin{Example}
    Suppose the scheduler determines the priority order MPC $>$ RRT $>$ TSP $>$ SLAM, and SLAM has the highest number of threads.
    Assigning SLAM a fixed priority value of 5, an acceptable mapping for the other tasks is MPC=8, RRT=7, TSP=6, SLAM=5.
    \end{Example}

    \subsection{Performance Metric Function}
    \label{section_generate_tsp_perf}
    Some tasks' output quality is either unaffected by hyper-parameters or has a good offline configuration, making online adjustment unnecessary; their performance metric function outputs a constant $1.0$.
    For other tasks, whose optimal configuration depends on the environment and available computation, online optimization is beneficial and requires a properly defined performance function.
    We use the TSP as an example.
    The TSP is an anytime algorithm, so longer running time limits yield shorter output paths.
    We profiled running time limits starting at 50 ms and increasing by 50 ms until further increases no longer shortened the path, testing each configuration 10 times and averaging.
    Where small increases did not improve the path length, the best result was carried forward, so the time limits in Table~\ref{table_tsp_perf} do not increase strictly in 50 ms steps.
    Path lengths were normalized to $[0, 1.0]$: the longest path scored 0.1 (still a usable solution; 0.0 is reserved for no valid output) and the shortest scored 1.0, with linear interpolation for intermediate values.

    \begin{table*}[]
    \centering
    \caption{TSP's performance metric function with different running time limits.}
    \begin{tabular}{@{}llllllllllllll@{}}
        \toprule
        Running Time Limit(ms)       & 150   & 250   & 300   & 450   & 850   & 900   & 950   & 1000  & 1050  & 1100  & 1200  & 1500   \\ \midrule
        Physical Distance  & 61551 & 52027 & 50969 & 50097 & 49759 & 49397 & 48915 & 48794 & 48527 & 48366 & 48244 & 48093 \\
        Performance Function & 0.1   & 0.73  & 0.80 & 0.86  & 0.88  & 0.91  & 0.94  & 0.95  & 0.97  & 0.98  & 0.99  & 1        \\ \bottomrule
    \end{tabular}
    \label{table_tsp_perf}
    \end{table*}

    Compared with the performance metric definition in Definition~\ref{def_perf_metric}, the TSP metric omits the ``environment'' factor for simplicity; environmental influences can be incorporated without affecting the optimization framework.

    \subsection{Hyperparameter Configuration}

    Only a few hyperparameters require configuration:
    \begin{itemize}
        \item The number of value-probability pairs describing a probability mass function during the fixed-point iteration in~\eqref{eq_prob_rta}.
        \item The number of partial priority assignments kept during priority assignment optimization ($m$ in Algorithm~\ref{alg:modified_audsley}, Section~\ref{section_audsley_pa}).
        \item The local search radius $\delta$ of the incremental priority move (Section~\ref{section_increment_pa}).
        \item The scheduler's period $T$ (Section~\ref{section_implementation}).
    \end{itemize}
    The first trades off computation cost against modeling accuracy; we use 10, beyond which no observable improvement justifies the extra cost.
    The second trades off computation cost against optimization quality; to keep overall overhead below 1\%, we set it to 2.
    For the third, we set $\delta$ so that only 3 candidates are considered each time.
    For the fourth, we use the lowest value keeping scheduler overhead below 1\%, namely 15~s in our real-world experiments.
